% c17 原书 p.210–223（PDF p.217–230）

\subsection{高阶直接像的局部自由性}\label{sec:2.6}

设 $f\colon X\to Y$ 是态射，$\cF$ 是拟凝聚 $\cO_X$-模。若对每个 $x\in X$，$\cF_x$ 是平坦 $\cO_{Y,f(x)}$-模，则称 $\cF$ \emph{在 $Y$ 上平坦}。由 \ref{it:2.1.19}，这等价于说：对 $X$ 的每个仿射开子概形 $U$ 以及 $Y$ 的每个满足 $f(U)\subset V$ 的仿射开子概形 $V$，$\cF(U)$ 是平坦 $\cO_Y(V)$-模。对每个 $y\in Y$，令 $k(y)$ 为局部环 $\cO_{Y,y}$ 的剩余域，令 $X_y=f^{-1}(y)=X\tensor_Y\Spec k(y)$，并令 $\cF_y=p^*\cF$，其中 $p\colon X\tensor_Y\Spec k(y)\to X$ 是投影。本节研究当 $y$ 遍历 $Y$ 时 $H^n(X_y,\cF_y)$ 如何变化。我们假设 $f$ 是诺特概形的固有态射，且 $\cF$ 是平坦于 $Y$ 的凝聚 $\cO_X$-模。此外，我们还得到 $R^if_*\cF$ 局部自由的一些判别法。

\begin{lemma}\label{it:2.6.1}
设 $A$ 是环，$K^\cdot$ 与 $C^\cdot$ 是平坦 $A$-模的两个复形，且当 $i$ 充分大时 $K^i=C^i=0$。设 $\phi\colon K^\cdot\to C^\cdot$ 是使上同调群同构的复形态射。则对每个 $A$-模 $M$，态射 $\phi\tensor\id_M\colon K^\cdot\tensor_A M\to C^\cdot\tensor_A M$ 都诱导上同调群的同构。
\end{lemma}

\begin{proof}
设 $L^\cdot$ 是态射 $\phi\colon K^\cdot\to C^\cdot$ 的映射锥。则 $L^\cdot\tensor_A M$ 是 $\phi\tensor\id_M$ 的映射锥。假设 $\phi$ 诱导上同调群的同构，则对每个 $i$ 都有 $H^i(L^\cdot)=0$。再假设当 $i>n$ 时 $K^i=C^i=0$。则当 $i>n$ 时 $L^i=0$。我们有短正合序列
\[
0\to Z^{n-1}(L^\cdot)\to L^{n-1}\to L^n\to 0,
\]
\[
0\to Z^{n-2}(L^\cdot)\to L^{n-2}\to Z^{n-1}(L^\cdot)\to 0,
\]
\[
0\to Z^{n-3}(L^\cdot)\to L^{n-3}\to Z^{n-2}(L^\cdot)\to 0,
\]
\[\vdots\]
若 $K^\cdot$ 与 $C^\cdot$ 是平坦 $A$-模的复形，则 $L^\cdot$ 也是平坦 $A$-模的复形。对任意 $A$-模 $M$，与上述短正合序列相伴的 $\Tor_i(-,M)$ 长正合序列依次表明
\[
\Tor_i(Z^{n-1}(L^\cdot),M)=0,\quad \Tor_i(M,Z^{n-2}(L^\cdot))=0,\quad\dots
\]
对所有 $i\ge 1$ 成立，且下列序列正合：
\[
0\to Z^{n-1}(L^\cdot)\tensor_A M\to L^{n-1}\tensor_A M\to L^n\tensor_A M\to 0,
\]
\[
0\to Z^{n-2}(L^\cdot)\tensor_A M\to L^{n-2}\tensor_A M\to Z^{n-1}(L^\cdot)\tensor_A M\to 0,
\]
\[
0\to Z^{n-3}(L^\cdot)\tensor_A M\to L^{n-3}\tensor_A M\to Z^{n-2}(L^\cdot)\tensor_A M\to 0,
\]
\[\vdots\]
所以
\[
\cdots\to L^i\tensor_A M\to L^{i+1}\tensor_A M\to\cdots
\]
正合。因此对任意 $i$ 都有 $H^i(L^\cdot\tensor_A M)=0$。这说明 $\phi\tensor\id_M\colon K^\cdot\tensor_A M\to C^\cdot\tensor_A M$ 诱导上同调群的同构。
\end{proof}

\begin{lemma}\label{it:2.6.2}
设 $A$ 是环，$K^\cdot$ 与 $C^\cdot$ 是平坦 $A$-模的两个复形，且当 $i$ 充分大时 $K^i=C^i=0$。设 $\phi\colon K^\cdot\to C^\cdot$ 是使上同调群同构的复形态射。假设对所有 $i<0$ 都有 $C^i=0$。则 $K^0/B^0(K^\cdot)$ 是平坦 $A$-模，且由 $\phi$ 诱导的复形态射
\[
\begin{tikzcd}[column sep=large]
0 \arrow[r] & K^0/B^0(K^\cdot) \arrow[r] \arrow[d] & K^1 \arrow[r] \arrow[d] & K^2 \arrow[r] \arrow[d] & \cdots \\
0 \arrow[r] & C^0 \arrow[r] & C^1 \arrow[r] & C^2 \arrow[r] & \cdots
\end{tikzcd}
\]
诱导上同调群的同构。
\end{lemma}

\begin{proof}
由于对所有 $i<0$ 都有 $C^i=0$ 且 $H^i(K^\cdot)\cong H^i(C^\cdot)$，我们对所有 $i<0$ 都有 $H^i(K^\cdot)=0$，从而
\[
\cdots\to K^{-1}\to K^0\to K^0/B^0(K^\cdot)\to 0
\]
是 $K^0/B^0(K^\cdot)$ 由平坦 $A$-模给出的分解。于是对每个 $A$-模 $M$，对所有 $i\ge 0$ 都有
\[
\Tor_i^A(K^0/B^0(K^\cdot),M)\cong H^{-i}(K^\cdot\tensor_A M).
\]
由 \ref{it:2.6.1}，我们有
\[
H^{-i}(K^\cdot\tensor_A M)\cong H^{-i}(C^\cdot\tensor_A M).
\]
由假设对所有 $i<0$ 都有 $C^i=0$，我们得到：对所有 $i\ge 1$ 都有
\[
H^{-i}(K^\cdot\tensor_A M)=0.
\]
所以对所有 $i\ge 1$ 与所有 $A$-模 $M$ 都有 $\Tor_i^A(K^0/B^0(K^\cdot),M)=0$。因此 $K^0/B^0(K^\cdot)$ 是平坦的。引理中其余的断言是显然的。
\end{proof}

\begin{lemma}\label{it:2.6.3}
设 $A$ 是诺特环，$C^\cdot$ 是平坦 $A$-模的复形，使得每个 $H^i(C^\cdot)$ 都是有限生成 $A$-模，且对所有 $i\notin[0,n]$ 都有 $C^i=0$。则存在有限生成投射 $A$-模的复形 $K^\cdot$ 与态射 $\phi\colon K^\cdot\to C^\cdot$，使得对所有 $i\notin[0,n]$ 都有 $K^i=0$，且对每个 $A$-模 $M$，$\phi\tensor\id_M\colon K^\cdot\tensor_A M\to C^\cdot\tensor_A M$ 诱导上同调群的同构。
\end{lemma}

\begin{proof}
对每个 $i>n$，令 $K^i=0$。设对每个 $i>p$，我们已经定义了 $K^i$ 与 $K^i\to C^i$，使得每个 $K^i$ 都是有限秩自由 $A$-模，$H^i(K^\cdot)\cong H^i(C^\cdot)$ 对每个 $i>p+1$ 成立，且 $Z^{p+1}(K^\cdot)\to H^{p+1}(C^\cdot)$ 是满射。考虑图表

\begin{center}
\begin{tikzcd}
A \arrow[r,"\alpha"] \arrow[d] & Z^{p+1}(K^\cdot) \arrow[d] \\
C^p \arrow[r] & Z^{p+1}(C^\cdot),
\end{tikzcd}
\end{center}

其中
\[
A=\{(c,k)\in C^p\oplus Z^{p+1}(K^\cdot)\mid c \text{ 与 } k \text{ 在 } Z^{p+1}(C^\cdot) \text{ 中有相同的像}\}.
\]
右侧竖直箭头诱导单态射 $\coker\alpha\to H^{p+1}(C^\cdot)$。由归纳假设，它是满射，从而是同构。图表中左侧竖直箭头诱导同构 $\ker\alpha\cong Z^p(C^\cdot)$。由假设，$H^p(C^\cdot)$ 是有限生成的。选取 $a_1,\dots,a_r\in\ker\alpha$，使得它们的像生成 $H^p(C^\cdot)$。由于 $K^{p+1}$ 是有限生成的且 $A$ 是诺特的，我们可以选取 $a_{r+1},\dots,a_s\in A$，使得 $\alpha(a_{r+1}),\dots,\alpha(a_s)$ 生成 $\im\alpha$。令 $K^p$ 是秩为 $s$ 的自由 $A$-模，并令 $\gamma\colon K^p\to A$ 是把 $K^p$ 的一组基映到 $\{a_1,\dots,a_r,a_{r+1},\dots,a_s\}$ 的同态。则 $\alpha\gamma$ 与 $\alpha$ 的像相同。这说明
\[
H^{p+1}(K^\cdot)\cong\coker(\alpha\gamma)\cong\coker\alpha\cong H^{p+1}(C^\cdot).
\]
由于 $\gamma$ 的像包含 $a_1,\dots,a_r$，$Z^p(K^\cdot)=\ker(\alpha\gamma)$ 满射到 $H^p(C^\cdot)$。这样我们得到有限秩自由 $A$-模的复形 $K^\cdot$ 与诱导上同调同构的态射 $K^\cdot\to C^\cdot$。此外，只有当 $i\le n$ 时才有 $K^i\ne 0$。为了得到具有所需性质的复形，我们把 $K^\cdot$ 换成
\[
0\to K^0/B^0(K^\cdot)\to K^1\to\cdots\to K^n\to 0.
\]
由 \ref{it:2.6.2}，$K^0/B^0(K^\cdot)$ 是平坦 $A$-模。由于 $K^0$ 与 $K^{-1}$ 都是有限秩自由 $A$-模，它具有有限表现。所以由 \ref{it:2.1.18}，$K^0/B^0(K^\cdot)$ 是投射 $A$-模。由 \ref{it:2.6.1} 与 \ref{it:2.6.2}，复形
\[
0\to K^0/B^0(K^\cdot)\to K^1\to\cdots\to K^n\to 0
\]
具有所需性质。
\end{proof}

\begin{proposition}\label{it:2.6.4}
设 $A$ 是诺特环，$Y=\Spec A$，$f\colon X\to Y$ 是固有态射，$\cF$ 是平坦于 $Y$ 的凝聚 $\cO_X$-模。则存在形如
\[
0\to K^0\to K^1\to\cdots\to K^n\to 0
\]
的复形 $K^\cdot$，使得对每个 $i$，$K^i$ 都是有限生成投射 $A$-模，且对每个 $A$-模 $M$，存在对 $M$ 自然同构
\[
H^i(X,\cF\tensor_A M)\cong H^i(K^\cdot\tensor_A M),
\]
其中 $\cF\tensor_A M=\cF\tensor_{\cO_X}f^*(M^\sim)$。此外，对每个 $A$-代数 $B$ 与每个 $i$，存在对 $B$ 自然同构
\[
H^i(X\tensor_Y\Spec B,\cF\tensor_A B)\cong H^i(K\tensor_A B),
\]
其中 $\cF\tensor_A B$ 是 $\cF$ 在投影 $p\colon X\tensor_Y\Spec B\to X$ 下的逆像 $p^*\cF$。特别地，对每个 $i$ 与每个 $y\in Y$，我们有
\[
H^i(X_y,\cF_y)\cong H^i(K^\cdot\tensor_A k(y))\cong H^i(X,\cF\tensor_A k(y)).
\]
\end{proposition}

\begin{proof}
设 $\fU=\{U_0,\dots,U_n\}$ 是 $X$ 的由有限多个仿射开子概形组成的覆盖。考虑交错 Čech 复形 $C^\cdot=C'^\cdot(\fU,\cF)$。我们有
\[
C^i=\prod_{k_0<\cdots<k_i}\cF(U_{k_0\dots k_i}),
\]
\[
C'^i(\fU,\cF\tensor_A M)=\prod_{k_0<\cdots<k_i}\cF(U_{k_0\dots k_i})\tensor_A M=C^i\tensor_A M.
\]
所以我们有 $C'^\cdot(\fU,\cF\tensor_A M)=C^\cdot\tensor_A M$。由 \ref{it:2.4.5}，对每个 $i$ 都有
\[
H^i(X,\cF\tensor_A M)\cong H^i(C^\cdot\tensor_A M).
\]
对任意 $A$-代数 $B$，$\fU'=\{p^{-1}(U_0),\dots,p^{-1}(U_n)\}$ 是 $X\tensor_Y\Spec B$ 的由仿射开子概形组成的覆盖。容易验证
\[
C'^\cdot(\fU',\cF\tensor_A B)=C^\cdot\tensor_A B.
\]
再次由 \ref{it:2.4.5}，对每个 $i$ 都有
\[
H^i(X\tensor_Y\Spec B,\cF\tensor_A B)\cong H^i(C^\cdot\tensor_A M).
\]
由于 $\cF$ 在 $Y$ 上平坦，$C^\cdot$ 是平坦 $A$-模的复形。由 \ref{it:2.5.3}，对每个 $i$，$H^i(C^\cdot)$ 都是有限生成 $A$-模。此外，只有当 $0\le i\le n$ 时才有 $C^i\ne 0$。于是我们的命题由 \ref{it:2.6.3} 得到。
\end{proof}

\begin{tnote}
命题 \ref{it:2.6.4} 证明中第二式右侧原书排印为 $H^i(C^\cdot\tensor_A M)$，按上下文应为 $H^i(C^\cdot\tensor_A B)$（$B$ 是任意 $A$-代数）；译文保留原书排印，读者可自行按 $B$ 理解。[依据：原书 p.214]
\end{tnote}

回忆：函数 $g\colon Y\to\bR$ 称为\emph{上半连续}\idx{shangbanlianxu}{上半连续}的，若对任意 $r\in\bR$，$g^{-1}((-\infty,r))$ 是开集。

\begin{theorem}[上半连续性定理]\label{it:2.6.5}
设 $f\colon X\to Y$ 是诺特概形的固有态射，$\cF$ 是平坦于 $Y$ 的凝聚 $\cO_X$-模。
\begin{enumerate}[label=(\roman*)]
\item 对每个 $i$，函数
\[
y\mapsto\dim_{k(y)}H^i(X_y,\cF_y)
\]
在 $Y$ 上上半连续。
\item 函数
\[
y\mapsto\chi(X_y,\cF_y)=\sum_i(-1)^i\dim_{k(y)}H^i(X_y,\cF_y)
\]
在 $Y$ 上局部常值\idx{jubuchangzhi}{局部常值}。
\end{enumerate}
\end{theorem}

\begin{proof}
问题关于 $Y$ 是局部的。于是我们可以假设 $Y=\Spec A$，其中 $A$ 是诺特环。设 $K^\cdot$ 是 \ref{it:2.6.4} 中的复形。我们有
\begin{align*}
\dim_{k(y)}H^i(X_y,\cF_y)
&=\dim_{k(y)}H^i(K^\cdot\tensor_A k(y))\\
&=\dim_{k(y)}\ker(d^i\tensor\id_{k(y)})-\dim_{k(y)}\im(d^{i-1}\tensor\id_{k(y)})\\
&=\dim_{k(y)}(K^i\tensor_A k(y))-\dim_{k(y)}\im(d^i\tensor\id_{k(y)})\\
&\qquad-\dim_{k(y)}\im(d^{i-1}\tensor\id_{k(y)}).
\end{align*}
把上述维数取交错和，我们得到
\[
\chi(X_y,\cF_y)=\sum_i(-1)^i\dim_{k(y)}(K^i\tensor_A k(y)).
\]
由于 $K^i$ 是有限生成投射 $A$-模，函数
\[
y\mapsto\dim_{k(y)}(K^i\tensor_A k(y))
\]
由 \ref{it:1.4.7} 在 $Y$ 上局部常值。所以函数
\[
y\mapsto\chi(X_y,\cF_y)=\sum_i(-1)^i\dim_{k(y)}H^i(X_y,\cF_y)
\]
在 $Y$ 上局部常值。这证明了 (ii)。为证明 (i)，只需证明对每个 $i$，函数
\[
y\mapsto\dim_{k(y)}\im(d^i\tensor\id_{k(y)})
\]
在 $Y$ 上下半连续。问题仍然是局部的。再次由 \ref{it:1.4.7}，我们可以假设 $K^i$ 与 $K^{i+1}$ 是有限秩自由 $A$-模。取定 $K^i$ 与 $K^{i+1}$ 的基，并令 $P$ 是同态 $d^i\colon K^i\to K^{i+1}$ 关于这些基的矩阵。我们有
\[
\{y\in Y\mid\dim_{k(y)}\im(d^i\tensor\id_{k(y)})\le r\}=V(\fa),
\]
其中 $\fa$ 是 $A$ 中由矩阵 $P$ 的所有 $(r+1)\times(r+1)$ 阶子式生成的理想。所以 $\{y\in Y\mid\dim_{k(y)}\im(d^i\tensor\id_{k(y)})\le r\}$ 是闭集，从而函数 $y\mapsto\dim_{k(y)}\im(d^i\tensor\id_{k(y)})$ 下半连续\idx{xiabanlianxu}{下半连续}。
\end{proof}

\begin{theorem}\label{it:2.6.6}
设 $f\colon X\to Y$ 是诺特概形的固有态射，$\cF$ 是平坦于 $Y$ 的凝聚 $\cO_X$-模。若 $Y$ 是既约的，且对某个 $i$，函数
\[
y\mapsto\dim_{k(y)}H^i(X_y,\cF_y)
\]
在 $Y$ 上局部常值，则 $R^if_*\cF$ 局部自由，且标准同态
\[
(R^if_*\cF)_y\tensor_{\cO_{Y,y}}k(y)\to H^i(X_y,\cF_y)
\]
对所有 $y\in Y$ 都是同构。
\end{theorem}

\begin{proof}
问题关于 $Y$ 是局部的。我们可以假设 $Y=\Spec A$，其中 $A$ 是既约诺特环。设 $K^\cdot$ 是 \ref{it:2.6.4} 中的复形。由假设，函数
\[
y\mapsto\dim_{k(y)}H^i(X_y,\cF_y)
\]
在 $Y$ 上局部常值。由于 $K^i$ 是投射的，函数
\[
y\mapsto\dim_{k(y)}(K^i\tensor_A k(y))
\]
由 \ref{it:1.4.7} 也局部常值。我们在 \ref{it:2.6.5} 的证明中已看到
\begin{align*}
&\dim_{k(y)}H^i(X_y,\cF_y)+\dim_{k(y)}\im(d^{i-1}\tensor\id_{k(y)})\\
&=\dim_{k(y)}(K^i\tensor_A k(y))-\dim_{k(y)}\im(d^i\tensor\id_{k(y)}).
\end{align*}
由左端定义的 $Y$ 上的函数是下半连续的，而由右端定义的函数是上半连续的。所以它们定义 $Y$ 上的连续函数。由于该函数取值于整数，它必局部常值。于是函数
\[
y\mapsto\dim_{k(y)}\im(d^i\tensor\id_{k(y)}),\qquad
y\mapsto\dim_{k(y)}\im(d^{i-1}\tensor\id_{k(y)})
\]
在 $Y$ 上局部常值。由于函数 $y\mapsto\dim_{k(y)}(K^{i+1}\tensor_A k(y))$ 由 \ref{it:1.4.7} 局部常值，函数
\[
y\mapsto\dim_{k(y)}\coker(d^i\tensor\id_{k(y)})
=\dim_{k(y)}(K^{i+1}\tensor_A k(y))-\dim_{k(y)}\im(d^i\tensor\id_{k(y)})
\]
也局部常值。由于函子 $-\tensor_A k(y)$ 右正合，我们有
\[
\coker(d^i\tensor\id_{k(y)})\cong(\coker d^i)\tensor_A k(y).
\]
所以函数 $y\mapsto\dim_{k(y)}((\coker d^i)\tensor_A k(y))$ 在 $Y$ 上局部常值。由下面的 \ref{it:2.6.7}，$\coker d^i$ 是投射 $A$-模。于是短正合序列
\[
0\to\im d^i\to K^{i+1}\to\coker d^i\to 0
\]
分裂，从而
\[
K^{i+1}\cong\im d^i\oplus\coker d^i.
\]
结合 $K^{i+1}$ 是投射的这一事实，我们看到 $\im d^i$ 是投射的。所以短正合序列
\[
0\to\ker d^i\to K^i\to\im d^i\to 0
\]
分裂。我们有分裂短正合序列
\[
0\to(\im d^i)\tensor_A k(y)\to K^{i+1}\tensor_A k(y)\to(\coker d^i)\tensor_A k(y)\to 0,
\]
\[
0\to(\ker d^i)\tensor_A k(y)\to K^i\tensor_A k(y)\to(\im d^i)\tensor_A k(y)\to 0.
\]
由此推出
\[
(\ker d^i)\tensor_A k(y)\cong\ker(d^i\tensor\id_{k(y)}).
\]
考虑交换图

\begin{center}
\begin{tikzcd}[column sep=small]
(\im d^{i-1})\tensor_A k(y) \arrow[r] \arrow[d] & (\ker d^i)\tensor_A k(y) \arrow[r] \arrow[d,"\cong"] & H^i(K^\cdot)\tensor_A k(y) \arrow[r] \arrow[d] & 0 \\
\im(d^{i-1}\tensor\id_{k(y)}) \arrow[r] & \ker(d^i\tensor\id_{k(y)}) \arrow[r] & H^i(K^\cdot\tensor_A k(y)) \arrow[r] & 0.
\end{tikzcd}
\end{center}

两行都正合。第一个竖直箭头是满射。中间的竖直箭头是同构。由此推出
\[
H^i(K^\cdot)\tensor_A k(y)\cong H^i(K^\cdot\tensor_A k(y)).
\]
特别地，函数
\[
y\mapsto\dim_{k(y)}(H^i(K^\cdot)\tensor_A k(y))
\]
局部常值。由下面的 \ref{it:2.6.7}，$(H^i(K^\cdot))^\sim$ 局部自由。由 \ref{it:2.6.4}，我们有
\[
H^i(X,\cF)\cong H^i(K^\cdot),\qquad H^i(X_y,\cF_y)\cong H^i(K^\cdot\tensor_A k(y)).
\]
由 \ref{it:2.4.9}，我们有 $R^if_*\cF\cong(H^i(X,\cF))^\sim$。所以
\[
R^if_*\cF\cong(H^i(K^\cdot))^\sim.
\]
因此 $R^if_*\cF$ 局部自由。此外，我们有
\[
H^i(X_y,\cF_y)\cong H^i(K^\cdot\tensor_A k(y))\cong H^i(K^\cdot)\tensor_A k(y)\cong(R^if_*\cF)_y\tensor_{\cO_{Y,y}}k(y).
\]
\end{proof}

\begin{lemma}\label{it:2.6.7}
设 $Y$ 是既约诺特概形，$\cF$ 是凝聚 $\cO_Y$-模。若函数
\[
y\mapsto\dim_{k(y)}(\cF_y\tensor_{\cO_{Y,y}}k(y))
\]
在 $Y$ 上局部常值，则 $\cF$ 局部自由。
\end{lemma}

\begin{proof}
设 $P$ 是 $Y$ 中的一点。选取 $P$ 的一个仿射开连通邻域 $V$ 以及截面 $s_1,\dots,s_n\in\cF(V)$，使得它们在 $\cF_P\tensor_{\cO_{Y,P}}k(P)$ 中的像是 $k(P)$-向量空间 $\cF_P\tensor_{\cO_{Y,P}}k(P)$ 的一组基。考虑态射
\[
\cO_V^n\to\cF|_V,\qquad e_i\mapsto s_i,
\]
其中对每个 $i$，$e_i$ 是 $\cO_V^n$ 的第 $i$ 个分量为 $1$、其余分量为 $0$ 的截面。令 $C$ 是该态射的余核。由中山引理，$(\cO_V^n)_P\to\cF_P$ 是满射，从而 $C_P=0$。由 \ref{it:1.4.1}，在缩小 $P$ 的邻域 $V$ 之后，我们可以假设 $C|_V=0$，从而 $\cO_V^n\to\cF|_V$ 是满射。由于 $V$ 连通，局部常值函数 $y\mapsto\dim_{k(y)}(\cF_y\tensor_{\cO_{Y,y}}k(y))$ 在 $V$ 上是常值。所以对每个 $y\in V$ 都有 $\dim_{k(y)}(\cF_y\tensor_{\cO_{Y,y}}k(y))=n$。结合 $\cO_V^n\to\cF|_V$ 是满射，我们看到
\[
(\cO_V^n)_y\tensor_{\cO_{Y,y}}k(y)\to\cF_y\tensor_{\cO_{Y,y}}k(y)
\]
对每个 $y\in V$ 都是同构。于是对 $\ker(\cO_V^n\to\cF)$ 在某个仿射开子集 $U\subset V$ 上的每个截面 $(a_1,\dots,a_n)$，每个 $a_i$ 在每点 $y\in U$ 处的芽都落在 $\cO_{Y,y}$ 的极大理想中。由于 $U$ 是既约的，我们必须有 $a_i=0$。所以满射 $\cO_V^n\to\cF$ 的核为 $0$。因此 $\cO_V^n\cong\cF$，且 $\cF$ 局部自由。
\end{proof}

\begin{proposition}\label{it:2.6.8}
设 $A$ 是极大理想为 $\fm$ 的诺特局部环，$X$ 是固有于 $\Spec A$ 的概形，$\cF$ 是平坦于 $\Spec A$ 的凝聚 $\cO_X$-模。对每个 $i$，下列条件等价：
\begin{enumerate}[label=(\roman*)]
\item 标准同态
\[
H^i(X,\cF)\tensor_A A/\fm\to H^i(X,\cF\tensor_A A/\fm)
\]
是满射。
\item 对每个 $A$-模 $M$，标准同态
\[
H^i(X,\cF)\tensor_A M\to H^i(X,\cF\tensor_A M)
\]
是满射。
\item $A$-模范畴上的函子
\[
M\mapsto H^i(X,\cF\tensor_A M)
\]
右正合。
\item 对每个 $A$-模 $M$，标准同态
\[
H^i(X,\cF)\tensor_A M\to H^i(X,\cF\tensor_A M)
\]
是同构。
\end{enumerate}
\end{proposition}

\begin{proof}
每个 $x\in M$ 定义态射
\[
\cF\to\cF\tensor_A M,\qquad s\mapsto s\tensor x,
\]
从而定义同态
\[
\phi_x\colon H^i(X,\cF)\to H^i(X,\cF\tensor_A M).
\]
标准同态
\[
H^i(X,\cF)\tensor_A M\to H^i(X,\cF\tensor_A M)
\]
由 $s\tensor x\mapsto\phi_x(s)$ 定义。
\begin{enumerate}[label=(\roman*)]
\item[(i)$\Rightarrow$(ii)] 我们先证明对每个有限长 $A$-模 $M$，
\[
H^i(X,\cF)\tensor_A M\to H^i(X,\cF\tensor_A M)
\]
是满射。当 $\operatorname{length}_AM=0$（即 $M=0$）时这显然。设 $\operatorname{length}_AM=1$。则对每个非零元 $x\in M$ 都有 $M=Ax$。于是 $M\cong A/\operatorname{Ann}(x)$。由于
\[
0\subset\fm/\operatorname{Ann}(x)\subset A/\operatorname{Ann}(x)
\]
是 $A$-模 $A/\operatorname{Ann}(x)$ 中的链，且 $\operatorname{length}_A(A/\operatorname{Ann}(x))=1$，我们必须有 $\fm=\operatorname{Ann}(x)$，从而 $M\cong A/\fm$。我们的断言于是由假设 (i) 得到。设我们已经对长度小于 $n$ 的 $A$-模证明了 (ii)，并设 $M$ 是长度为 $n$ 的 $A$-模。则我们可以找到 $A$-模的短正合序列
\[
0\to M'\to M\to M''\to 0,
\]
使得 $\operatorname{length}_A(M'),\operatorname{length}_A(M'')<n$。考虑交换图

\begin{center}
\begin{tikzcd}[column sep=small]
H^i(X,\cF)\tensor_A M' \arrow[r] \arrow[d] & H^i(X,\cF)\tensor_A M \arrow[r] \arrow[d] & H^i(X,\cF)\tensor_A M'' \arrow[r] \arrow[d] & 0 \\
H^i(X,\cF\tensor_A M') \arrow[r] & H^i(X,\cF\tensor_A M) \arrow[r] & H^i(X,\cF\tensor_A M'').
\end{tikzcd}
\end{center}

由于函子 $H^i(X,\cF)\tensor_A-$ 右正合，第一行正合。由于 $\cF$ 平坦于 $\Spec A$，序列
\[
0\to\cF\tensor_A M'\to\cF\tensor_A M\to\cF\tensor_A M''\to 0
\]
正合。所以上图中的第二行也正合。由归纳假设，第一与第三个竖直箭头是满射。这说明第二个竖直箭头是满射，从而 (ii) 对每个有限长 $A$-模成立。

其次我们证明 (ii) 对每个有限生成 $A$-模 $M$ 成立。注意对每个 $n$，$M/\fm^nM$ 是有限长 $A$-模。所以
\[
H^i(X,\cF)\tensor_A M/\fm^nM\to H^i(X,\cF\tensor_A M/\fm^nM)
\]
是满射。设 $K_n$ 是它的核。每个 $K_n$ 都有有限长。于是逆向系统 $(K_n)_{n\in\bN}$ 满足 Mittag-Leffler 条件。因此由 \ref{it:1.5.1}，同态
\[
\plim_n(H^i(X,\cF)\tensor_A M/\fm^nM)\to\plim_nH^i(X,\cF\tensor_A M/\fm^nM)
\]
是满射。令 $\widehat A$ 是 $A$ 的 $\fm$-进完全化。由 \ref{it:1.5.5}，我们有
\[
\plim_n(H^i(X,\cF)\tensor_A M/\fm^nM)\cong(H^i(X,\cF)\tensor_A M)^\wedge\cong(H^i(X,\cF)\tensor_A M)\tensor_A\widehat A.
\]
由 \ref{it:2.5.6}，我们有
\[
\plim_nH^i(X,\cF\tensor_A M/\fm^nM)\cong(H^i(X,\cF\tensor_A M))^\wedge\cong H^i(X,\cF\tensor_A M)\tensor_A\widehat A.
\]
所以同态
\[
(H^i(X,\cF)\tensor_A M)\tensor_A\widehat A\to H^i(X,\cF\tensor_A M)\tensor_A\widehat A
\]
是满射。由 \ref{it:2.1.23}，$\widehat A$ 是忠实平坦 $A$-代数。所以
\[
H^i(X,\cF)\tensor_A M\to H^i(X,\cF\tensor_A M)
\]
是满射。这证明了对每个有限生成 $A$-模，(ii) 成立。

对每个 $A$-模 $M$，我们有 $M=\dlim_nM_n$，其中 $(M_n)$ 是 $M$ 的有限生成 $A$-子模构成的有向系统。我们有
\[
H^i(X,\cF)\tensor_A M\cong\dlim_n(H^i(X,\cF)\tensor_A M_n),
\]
\[
H^i(X,\cF\tensor_A M)\cong\dlim_nH^i(X,\cF\tensor_A M_n),
\]
其中第二个等式由 \ref{it:2.1.12} (ii) 得到。由于每个
\[
H^i(X,\cF)\tensor_A M_n\to H^i(X,\cF\tensor_A M_n)
\]
都是满射，
\[
H^i(X,\cF)\tensor_A M\to H^i(X,\cF\tensor_A M)
\]
也是满射。

\item[(ii)$\Rightarrow$(iii)] 给定 $A$-模的短正合序列
\[
0\to M'\to M\to M''\to 0,
\]
我们有交换图

\begin{center}
\begin{tikzcd}[column sep=small]
H^i(X,\cF)\tensor_A M' \arrow[r] \arrow[d] & H^i(X,\cF)\tensor_A M \arrow[r] \arrow[d] & H^i(X,\cF)\tensor_A M'' \arrow[r] \arrow[d] & 0 \\
H^i(X,\cF\tensor_A M') \arrow[r] & H^i(X,\cF\tensor_A M) \arrow[r] & H^i(X,\cF\tensor_A M''),
\end{tikzcd}
\end{center}

其中两行正合，且所有竖直箭头都是满射。这说明 $H^i(X,\cF\tensor_A M)\to H^i(X,\cF\tensor_A M'')$ 是满射。所以序列
\[
H^i(X,\cF\tensor_A M')\to H^i(X,\cF\tensor_A M)\to H^i(X,\cF\tensor_A M'')\to 0
\]
正合，从而函子 $H^i(X,\cF\tensor_A-)$ 右正合。

\item[(iii)$\Rightarrow$(iv)] 先假设 $M$ 是有限生成 $A$-模。则我们可以找到正合序列
\[
A^m\to A^n\to M\to 0,
\]
其中 $m,n$ 是自然数。考虑如下的交换图：

\begin{center}
\begin{tikzcd}[column sep=small]
H^i(X,\cF)\tensor_A A^m \arrow[r] \arrow[d] & H^i(X,\cF)\tensor_A A^n \arrow[r] \arrow[d] & H^i(X,\cF)\tensor_A M \arrow[r] \arrow[d] & 0 \\
H^i(X,\cF\tensor_A A^m) \arrow[r] & H^i(X,\cF\tensor_A A^n) \arrow[r] & H^i(X,\cF\tensor_A M) \arrow[r] & 0.
\end{tikzcd}
\end{center}

由于函子 $H^i(X,\cF)\tensor_A-$ 与 $H^i(X,\cF\tensor_A-)$ 都右正合，两行正合。前两个竖直箭头是同构。所以最后一个竖直箭头也是同构。这就对有限生成 $A$-模证明了 (iv)。一般情形通过把模写成其有限生成子模的有向极限得到。

\item[(iv)$\Rightarrow$(i)] 显然。
\end{enumerate}
\end{proof}

\begin{tnote}
命题 \ref{it:2.6.8} 的 (ii)$\Rightarrow$(iii) 一段中，原书排印为「$H^i(X,\cF\tensor M)\to H^i(X,\cF\tensor_A M')$ 是满射」，所指应为第二行的第二个箭头，即 $H^i(X,\cF\tensor_A M)\to H^i(X,\cF\tensor_A M'')$；译文按上下文写作 $M''$。另原书此段 $M'$、$M''$ 与中段的 $M$ 排印混用，读者按短正合序列 $0\to M'\to M\to M''\to 0$ 理解即可。[依据：原书 p.220]
\end{tnote}

\begin{theorem}\label{it:2.6.9}
设 $f\colon X\to Y$ 是诺特概形的固有态射，$\cF$ 是平坦于 $Y$ 的凝聚 $\cO_X$-模。假设对某个 $i$，标准同态
\[
(R^if_*\cF)_y\tensor_{\cO_{Y,y}}k(y)\to H^i(X_y,\cF_y)
\]
对每个 $y\in Y$ 都是满射。则它是同构。此外，在此假设下，下列两个条件等价：
\begin{enumerate}[label=(\roman*)]
\item $\cO_Y$-模 $R^if_*\cF$ 局部自由。
\item 标准同态
\[
(R^{i-1}f_*\cF)_y\tensor_{\cO_{Y,y}}k(y)\to H^{i-1}(X_y,\cF_y)
\]
对每个 $y\in Y$ 都是满射。
\end{enumerate}
\end{theorem}

\begin{proof}
利用 \ref{it:2.4.9} 与 \ref{it:2.4.10}，可以验证
\[
(R^if_*\cF)_y\cong H^i(X\tensor_Y\Spec\cO_{Y,y},p^*\cF),
\]
其中 $p\colon X\tensor_Y\Spec\cO_{Y,y}\to X$ 是投影。注意 $\cF_y$ 作为 $X_y$ 上的层在闭浸入 $X_y\to X\tensor_Y\Spec\cO_{Y,y}$ 下的直接像是 $X\tensor_Y\Spec\cO_{Y,y}$ 上的层 $p^*\cF\tensor_{\cO_{Y,y}}k(y)$。所以我们有
\[
H^i(X_y,\cF_y)\cong H^i(X\tensor_Y\Spec\cO_{Y,y},p^*\cF\tensor_{\cO_{Y,y}}k(y)).
\]
由假设，
\[
(R^if_*\cF)_y\tensor_{\cO_{Y,y}}k(y)\to H^i(X_y,\cF_y)
\]
是满射。所以标准同态
\[
H^i(X\tensor_Y\Spec\cO_{Y,y},p^*\cF)\tensor_{\cO_{Y,y}}k(y)\to H^i(X\tensor_Y\Spec\cO_{Y,y},p^*\cF\tensor_{\cO_{Y,y}}k(y))
\]
是满射。对固有于 $\Spec\cO_{Y,y}$ 的概形 $X\tensor_Y\Spec\cO_{Y,y}$ 与平坦于 $\Spec\cO_{Y,y}$ 的凝聚层 $p^*\cF$ 应用 \ref{it:2.6.8}，我们看到该同态是同构。这证明了定理的第一个断言。由 \ref{it:2.6.8}，我们对每个 $\cO_{Y,y}$-模 $M$ 都有
\[
H^i(X\tensor_Y\Spec\cO_{Y,y},p^*\cF)\tensor_{\cO_{Y,y}}M\cong H^i(X\tensor_Y\Spec\cO_{Y,y},p^*\cF\tensor_{\cO_{Y,y}}M)
\]
由 \ref{it:1.4.1} (ii)，$\cO_Y$-模 $R^if_*\cF$ 局部自由，当且仅当对每个 $y\in Y$，$(R^if_*\cF)_y\cong H^i(X\tensor_Y\Spec\cO_{Y,y},p^*\cF)$ 是自由 $\cO_{Y,y}$-模。由 \ref{it:2.1.18} 与 \ref{it:1.4.8}，这等价于说 $\cO_{Y,y}$-模范畴上的函子
\[
M\mapsto H^i(X\tensor_Y\Spec\cO_{Y,y},p^*\cF)\tensor_{\cO_{Y,y}}M\cong H^i(X\tensor_Y\Spec\cO_{Y,y},p^*\cF\tensor_{\cO_{Y,y}}M)
\]
对每个 $y\in Y$ 都正合。利用 $\cF$ 平坦于 $Y$ 的假设以及上同调群的长正合序列，我们看到函子
\[
M\mapsto H^i(X\tensor_Y\Spec\cO_{Y,y},p^*\cF\tensor_{\cO_{Y,y}}M)
\]
正合，当且仅当下列两个函子右正合：
\[
M\mapsto H^i(X\tensor_Y\Spec\cO_{Y,y},p^*\cF\tensor_{\cO_{Y,y}}M),
\]
\[
M\mapsto H^{i-1}(X\tensor_Y\Spec\cO_{Y,y},p^*\cF\tensor_{\cO_{Y,y}}M).
\]
由 \ref{it:2.6.8}，这等价于说下列标准同态都是满射：
\[
H^i(X\tensor_Y\Spec\cO_{Y,y},p^*\cF)\tensor_{\cO_{Y,y}}k(y)\to H^i(X\tensor_Y\Spec\cO_{Y,y},p^*\cF\tensor_{\cO_{Y,y}}k(y)),
\]
\[
H^{i-1}(X\tensor_Y\Spec\cO_{Y,y},p^*\cF)\tensor_{\cO_{Y,y}}k(y)\to H^{i-1}(X\tensor_Y\Spec\cO_{Y,y},p^*\cF\tensor_{\cO_{Y,y}}k(y)),
\]
也就是说，
\[
(R^if_*\cF)_y\tensor_{\cO_{Y,y}}k(y)\to H^i(X_y,\cF_y),
\]
\[
(R^{i-1}f_*\cF)_y\tensor_{\cO_{Y,y}}k(y)\to H^{i-1}(X_y,\cF_y)
\]
都是满射。所以，在
\[
(R^if_*\cF)_y\tensor_{\cO_{Y,y}}k(y)\to H^i(X_y,\cF_y)
\]
对每个 $y\in Y$ 都是满射的假设下，$\cO_Y$-模 $R^if_*\cF$ 局部自由，当且仅当
\[
(R^{i-1}f_*\cF)_y\tensor_{\cO_{Y,y}}k(y)\to H^{i-1}(X_y,\cF_y)
\]
对每个 $y\in Y$ 都是满射。
\end{proof}

\begin{corollary}\label{it:2.6.10}
设 $f\colon X\to Y$ 是诺特概形的固有态射，$\cF$ 是平坦于 $Y$ 的 $\cO_X$-模。
\begin{enumerate}[label=(\roman*)]
\item 若对某个 $i$ 与每个 $y\in Y$ 都有 $H^i(X_y,\cF_y)=0$，则 $R^if_*\cF=0$ 且
\[
(R^{i-1}f_*\cF)_y\tensor_{\cO_{Y,y}}k(y)\cong H^{i-1}(X_y,\cF_y).
\]
\item 若对每个 $y\in Y$ 都有 $H^1(X_y,\cF_y)=0$，则 $\cO_Y$-模 $f_*\cF$ 局部自由。
\item 设 $i_0$ 是自然数，使得对所有 $i\ge i_0$ 都有 $R^if_*\cF=0$。则对所有 $i\ge i_0$ 与每个 $y\in Y$ 都有 $H^i(X_y,\cF_y)=0$。
\end{enumerate}
\end{corollary}

\begin{proof}
(i) 标准同态
\[
(R^if_*\cF)_y\tensor_{\cO_{Y,y}}k(y)\to H^i(X_y,\cF_y)
\]
对每个 $y\in Y$ 都平凡地是满射。由 \ref{it:2.6.9}，它是同构。所以对每个 $y\in Y$ 都有 $(R^if_*\cF)_y\tensor_{\cO_{Y,y}}k(y)=0$。由中山引理，对每个 $y\in Y$ 都有 $(R^if_*\cF)_y=0$，从而 $R^if_*\cF=0$。特别地，$R^if_*\cF=0$ 局部自由。于是由 \ref{it:2.6.9}，标准同态
\[
(R^{i-1}f_*\cF)_y\tensor_{\cO_{Y,y}}k(y)\to H^{i-1}(X_y,\cF_y)
\]
对每个 $y\in Y$ 都是同构。

(ii) 由 (i)，标准同态
\[
(R^0f_*\cF)_y\tensor_{\cO_{Y,y}}k(y)\to H^0(X_y,\cF_y)
\]
是满射，且标准同态
\[
(R^{-1}f_*\cF)_y\tensor_{\cO_{Y,y}}k(y)\to H^{-1}(X_y,\cF_y)
\]
对每个 $y\in Y$ 都平凡地是满射。所以由 \ref{it:2.6.9}，$f_*\cF=R^0f_*\cF$ 局部自由。

(iii) 我们知道对每个 $y\in Y$ 与充分大的 $i$ 都有 $H^i(X_y,\cF_y)=0$。假设 $i>i_0$ 且对每个 $y\in Y$ 都有 $H^i(X_y,\cF_y)=0$。则由 (i)，我们有
\[
H^{i-1}(X_y,\cF_y)\cong(R^{i-1}f_*\cF)_y\tensor_{\cO_{Y,y}}k(y)=0.
\]
\end{proof}

\begin{tnote}
推论 \ref{it:2.6.10} 的证明中 (ii) 用到 $R^{-1}f_*\cF$ 这一记号，即取 $-1$ 阶直接像为 $0$，故相应同态平凡地满射；原书未单独说明，此处按惯例保留。[依据：原书 p.222–223]
\end{tnote}
